\chapter{Integraciones, seguridad y observabilidad en detalle}
\label{anexo:integraciones}

Este anexo amplía las integraciones externas, las capas de seguridad y el stack de observabilidad resumidos en el Capítulo~\ref{cap:implementacion}.

% ============================================================
\section{Integración con Microsoft SharePoint}
\label{anexo:sharepoint}
% ============================================================

El módulo Indexer funciona como un servicio autónomo (\textit{worker}) que ejecuta un ciclo de sincronización periódica cada 5 minutos.

\paragraph{Autenticación MSAL.} La conexión con SharePoint se realiza mediante OAuth 2.0 con flujo de credenciales de cliente (\textit{Client Credentials Flow}). El servicio obtiene un token de acceso del \textit{tenant} de Azure AD mediante la librería MSAL y utiliza la API REST de Microsoft Graph \cite{msgraph2024docs} para listar y descargar ficheros.

\paragraph{Sincronización delta incremental.} Para evitar reprocesar documentos ya ingeridos, el Indexer consulta la marca temporal de modificación de cada fichero en SharePoint y la compara con los registros almacenados en PostgreSQL. Solo los ficheros nuevos o modificados se descargan, procesan e indexan. Esta estrategia reduce drásticamente el consumo de red y GPU en ciclos de sincronización recurrentes.

% ============================================================
\section{Resolución Semántica de Leyes (BOE)}
\label{anexo:boe}
% ============================================================

Para consultas como ``Resume la Ley de Protección de Datos'', el sistema implementa un mecanismo que mapea nombres coloquiales a identificadores oficiales del BOE:

\begin{lstlisting}[language=Python, caption={Mapeo semántico de nombres de leyes.}]
LAW_MAPPINGS = {
    "proteccion de datos": "BOE-A-2018-16673",
    "constitucion": "BOE-A-1978-31229",
    "propiedad intelectual": "BOE-A-1996-8930",
    "igualdad": "BOE-A-2007-6115",
}
\end{lstlisting}

El resolvedor elimina palabras vacías comunes del español (``la'', ``de'', ``del'', ``ley'') y busca coincidencias parciales en el diccionario antes de acudir al endpoint de búsqueda general del BOE.

En el alcance actual, el servidor MCP se ha validado en ejecución por terminal con un LLM y un cliente MCP compatible, pero la funcionalidad BOE \emph{dentro de la interfaz web} se resuelve mediante llamadas a la API REST del backend desde el pipeline JARVIS, no por MCP. Conviene explicar por qué, ya que es una fuente habitual de confusión.

\paragraph{Por qué el flujo web no pasa por MCP.} MCP es un protocolo cliente-servidor a nivel de aplicación: el cliente reside en la herramienta (Claude Desktop, OpenWebUI), no en el motor de inferencia. Ollama, por tanto, no «usa» MCP; se limita a servir el modelo. Para que OpenWebUI invocara las herramientas del servidor BOE por MCP tendrían que cumplirse dos condiciones que el diseño actual no satisface:

\begin{itemize}
    \item \textbf{Un modelo con \textit{function calling}.} El \textit{tool-calling} por MCP requiere que sea el propio LLM quien decida invocar la herramienta y construya sus argumentos. El modelo principal (\texttt{rag-qwen-ft}, alias JARVIS) está ajustado para responder en modo RAG con citas y se expone en LiteLLM con \texttt{supports\_function\_calling: false}; no emite llamadas a herramientas de forma fiable.
    \item \textbf{Que el modelo conduzca la conversación.} El \textit{tool-calling} de OpenWebUI opera en el diálogo modelo$\leftrightarrow$interfaz, pero en este sistema el método \texttt{pipe()} del pipeline JARVIS intercepta cada mensaje \emph{antes} y enruta por su cuenta (detección de intención por expresiones regulares). Como el pipeline ya decide todo, la capa de herramientas MCP de OpenWebUI nunca llega a activarse.
\end{itemize}

A esto se añade la cronología: el soporte MCP nativo de OpenWebUI maduró después de construir el pipeline, por lo que la vía integrada y validada fue API-vía-pipeline, que además da control directo sobre autenticación, filtrado por \textit{tenant}, métricas y formato de citas.

\paragraph{¿Merece la pena habilitar MCP de forma activa?} Para el caso concreto del BOE ---un dominio acotado, con siete herramientas y la resolución semántica ya resuelta en el pipeline--- la integración por API es suficiente y más controlada; activar MCP no aportaría una mejora proporcional a su coste: obligaría a adoptar un modelo con \textit{function calling} fiable, a reimplementar en la capa MCP las garantías de acceso y trazabilidad que hoy ofrece el pipeline, y a asumir mayor latencia y el riesgo de que el modelo elija o parametrice mal la herramienta (frente al determinismo del enrutamiento por reglas). El MCP activo sí compensaría en otro escenario: cuando se quieran incorporar \emph{muchas} herramientas o fuentes nuevas, interoperar con clientes MCP externos, o adoptar un modelo con \textit{tool-calling} robusto como núcleo del sistema; en ese caso MCP reduce el acoplamiento y escala mejor que seguir añadiendo ramas al pipeline. El camino intermedio más razonable es mantener el pipeline como orquestador pero hacer que sea \emph{él} quien actúe como cliente MCP de las herramientas nuevas, combinando el estándar con el control actual.

% ============================================================
\section{Validación JWT en el Backend}
\label{anexo:jwt}
% ============================================================

Además de la autenticación perimetral de OpenWebUI vía OIDC, el Backend implementa una segunda capa de validación directa de tokens JWT para clientes externos y herramientas MCP que llaman a los endpoints sin pasar por la interfaz web. La lógica, implementada en \texttt{backend/app/core/auth.py}, sigue el flujo estándar de validación de tokens Bearer de Azure AD:

\begin{enumerate}
    \item El backend descarga dinámicamente el conjunto de claves públicas del \textit{tenant} corporativo desde el endpoint JWKS de Azure AD.
    \item Las claves se cachean en memoria para evitar latencia en peticiones sucesivas.
    \item Cada token JWT recibido en la cabecera \texttt{Authorization: Bearer} se verifica criptográficamente contra esas claves.
    \item Los grupos del usuario extraídos del token se mapean a colecciones de Qdrant autorizadas.
    \item Si el token no es válido o el usuario carece de acceso a alguna colección, el backend devuelve \texttt{HTTP 403 Forbidden} antes de ejecutar ninguna búsqueda.
\end{enumerate}

Esta funcionalidad se activa con la variable de entorno \texttt{AZURE\_JWT\_VALIDATION=true} y es completamente transparente para los clientes que ya utilizan el flujo SSO de OpenWebUI.

\paragraph{Aislamiento de red.} Toda la malla de servicios convive dentro de un plano interno \texttt{rag-network}, una red Docker de tipo \textit{bridge}. Servicios como PostgreSQL (5432), Redis (6379), Qdrant (6333) y Ollama (11434) no son accesibles desde fuera del servidor bajo ninguna circunstancia, eliminando vectores de ataque sobre estas bases de datos.

% ============================================================
\section{Métricas, telemetría GPU y dashboards}
\label{anexo:observabilidad}
% ============================================================

\paragraph{Métricas personalizadas en el backend.} El backend FastAPI expone métricas en formato Prometheus a través del endpoint \texttt{/metrics}: contador de peticiones HTTP (por endpoint, método y código), histograma de latencia (para calcular p50, p95, p99), distribución de modos de operación (RAG, chat, web, scraping, OCR, BOE) y tasa de errores por tipo.

\paragraph{Telemetría GPU con NVIDIA DCGM Exporter.} Captura consumo de VRAM, utilización de CUDA cores, temperatura del chip y potencia consumida. Estas métricas resultaron críticas para determinar qué combinaciones de modelos podían coexistir en la VRAM disponible y cuándo recurrir a cuantizaciones más agresivas (Q4 en lugar de Q8).

\paragraph{Dashboards en Grafana.} Se diseñaron tres dashboards interactivos:
\begin{enumerate}
    \item Dashboard de Backend: latencia media y percentiles, throughput, tasa de errores, distribución de modos y tendencia de uso.
    \item Dashboard de GPU: consumo de VRAM con alarmas de umbral (80\%), temperatura, utilización de cores y potencia.
    \item Dashboard de Base de Datos: métricas de PostgreSQL (conexiones, queries/s) y de Qdrant (puntos indexados, latencia de búsqueda).
\end{enumerate}

\paragraph{Logs centralizados: Loki y Promtail.} Junto a las métricas, el sistema agrega los logs de todos los contenedores con \textbf{Promtail}, que los recolecta y los envía a \textbf{Loki}. Loki los indexa y los hace consultables desde el mismo Grafana, de modo que métricas y logs conviven en una única herramienta. Esto permite correlacionar un pico de latencia detectado en Prometheus con las líneas de log exactas de ese instante, acelerando el diagnóstico. La configuración reside en \texttt{config/promtail/config.yaml} y en el datasource \texttt{config/grafana/provisioning/datasources/loki.yml}.

\paragraph{Impacto de la monitorización.} El stack de observabilidad condujo a decisiones concretas: se validó que $\sim$35\% de las consultas repetidas se beneficiaban de la caché Redis; se detectó que las consultas RAG con más de 5 fragmentos no mejoraban la calidad pero aumentaban la latencia (optimizando $k$); y se observó que PaddleOCR monopolizaba la GPU durante la ingesta, lo que motivó programarla en horarios de baja actividad.
